% Responsable: Anthony · Adicional de la rúbrica · Ronda 4 (domingo 15, hasta 15.00)
% Rúbrica C3: actores del proceso (ingeniero de requisitos, cliente, usuario, desarrollador) con roles concretos en el restaurante, y actividades centrales y transversales del framework de Pohl mapeadas al sistema.
\subsection{Actores del proceso y framework de IR aplicado}
\subsubsection{Actores del proceso de requisitos}

El proceso de Ingeniería de Requisitos requiere la participación de actores
con responsabilidades diferentes. Para el PFC se identifican los roles de
ingeniero de requisitos, cliente, usuario y desarrollador, relacionándolos
con el contexto concreto de El Balcón Restaurant
\cite{ieee29148,pohl2025}.

\begin{table}[H]
\centering
\footnotesize
\caption{Actores del proceso de requisitos aplicados al PFC}
\label{tab:actores-ir}
\begin{tabularx}{\textwidth}{@{}p{0.21\textwidth}p{0.28\textwidth}Y@{}}
\toprule
\textbf{Actor del proceso} &
\textbf{Representación en el PFC} &
\textbf{Responsabilidad} \\
\midrule

Ingeniero de requisitos &
Equipo ACNN &
Elicitar, analizar, documentar, negociar y validar las necesidades del
restaurante, manteniendo coherencia y trazabilidad entre los artefactos. \\

Cliente &
Propietario de El Balcón Restaurant &
Expresar objetivos y necesidades del negocio, establecer prioridades y
aprobar los requisitos y el alcance acordado. \\

Usuario &
Jefe de cocina, meseros, cajero y clientes &
Aportar información sobre los procesos reales, las tareas cotidianas, las
dificultades y las condiciones de uso que debe considerar el sistema. \\

Desarrollador &
Equipo técnico encargado de implementar el PFC &
Transformar los requisitos aprobados en componentes de software, mantener
trazabilidad con la especificación e informar restricciones técnicas que
puedan afectar la solución. \\

\bottomrule
\end{tabularx}
\end{table}

\subsubsection{Actividades centrales del framework de Pohl}

Pohl organiza las actividades centrales de Ingeniería de Requisitos en
elicitación, documentación y negociación \cite{pohl2025}. Estas actividades
se relacionan entre sí y pueden repetirse durante el proyecto cuando se
descubre nueva información o se modifican las necesidades.

\begin{table}[H]
\centering
\footnotesize
\caption{Actividades centrales de IR aplicadas a El Balcón}
\label{tab:actividades-centrales}
\begin{tabularx}{\textwidth}{@{}p{0.20\textwidth}YY@{}}
\toprule
\textbf{Actividad} &
\textbf{Aplicación en el restaurante} &
\textbf{Artefacto o resultado} \\
\midrule

Elicitación &
Obtener información del propietario y los usuarios acerca del registro de
pedidos, reservas, inventario, ventas y necesidades de apoyo mediante IA. &
Stakeholders, contexto, preguntas de análisis y candidatos a requisitos. \\

Documentación &
Representar de forma estructurada las necesidades obtenidas para que puedan
ser revisadas y utilizadas durante el desarrollo. &
Tabla de requisitos, descripción del sistema, límite, visión y perspectivas
de contexto. \\

Negociación &
Resolver diferencias de prioridad, alcance o interpretación entre el
propietario, usuarios y equipo del proyecto. &
Requisitos acordados, prioridades y decisiones de alcance. \\

\bottomrule
\end{tabularx}
\end{table}

\subsubsection{Actividades transversales}

La validación y la gestión son actividades transversales porque acompañan a
las actividades centrales y a los artefactos producidos durante el proceso
\cite{pohl2025}.

\begin{table}[H]
\centering
\footnotesize
\caption{Actividades transversales aplicadas al PFC}
\label{tab:actividades-transversales}
\begin{tabularx}{\textwidth}{@{}p{0.20\textwidth}YY@{}}
\toprule
\textbf{Actividad} &
\textbf{Aplicación en El Balcón} &
\textbf{Resultado esperado} \\
\midrule

Validación &
Revisar con el cliente y los usuarios que los requisitos representen las
necesidades reales y evaluar que sean claros, verificables y coherentes. &
Corrección de ambigüedades, conflictos, omisiones e inconsistencias antes de
la implementación. \\

Gestión &
Controlar cambios, prioridades, versiones y relaciones de trazabilidad de
los requisitos durante el proyecto. &
Historial de cambios y correspondencia entre necesidades, requisitos,
artefactos y futuras pruebas. \\

\bottomrule
\end{tabularx}
\end{table}

\subsubsection{Mapeo integrado del framework}

La aplicación del framework permite relacionar los actores con las actividades
que realizan durante el proyecto. El propietario participa principalmente
como fuente de necesidades, negociador de prioridades y responsable de la
aceptación del alcance. Los usuarios aportan conocimiento operativo y
participan en la validación de los flujos de trabajo. El equipo ACNN organiza
la elicitación, documentación, negociación, validación y gestión de los
artefactos, mientras que el desarrollador utiliza los requisitos aceptados
como entrada para construir y verificar la solución.

Este mapeo demuestra que la Ingeniería de Requisitos del PFC no se limita a
redactar requisitos, sino que establece responsabilidades concretas,
actividades y artefactos para transformar las necesidades de El Balcón
Restaurant en una especificación controlada y trazable.
